% Copyright (c) 2026 Eugene Joseph M. Ragasa.
% SPDX-License-Identifier: Apache-2.0
% Author: Eugene Joseph M. Ragasa
% Formal proof completed: 2026-10-19
\documentclass[11pt]{article}
\usepackage[T1]{fontenc}
\usepackage{amsmath,amssymb,booktabs,tabularx}
\usepackage[margin=1in]{geometry}
\usepackage[hidelinks]{hyperref}
\title{Proof Dossier: Exact Dominance Check}
\author{Eugene Joseph M. Ragasa}
\date{Formal proof completed: 2026-10-19}
\begin{document}\maketitle
\section{Identity and status}
\begin{tabularx}{\textwidth}{@{}l>{\raggedright\arraybackslash}X@{}}\toprule
Statement ID & \texttt{BG-CERT-THM-010} \\
Status & proved; no admissions \\
Lean declaration & \path{ProjectKoios.BandGauge.Certificates.checkDominates_eq_true_iff} \\
Owning module & \path{ProjectKoios/BandGauge/Certificates/FinitePareto.lean} \\
Axioms & \texttt{propext}, \texttt{Classical.choice}, \texttt{Quot.sound} \\
\bottomrule\end{tabularx}
\section{Statement}
For a finite objective type with an exact decidable linear order, the Boolean dominance
check returns true exactly when the first decision weakly improves the second in every
objective and strictly improves it in at least one objective.
\section{Argument}
The checker filters the complete objective enumeration once for weak inequalities and
once for strict inequalities. The first filter must equal the full enumeration and the
second must be nonempty. Simplifying finite-filter membership gives exactly the two
conjuncts in the propositional dominance definition.
\section{Boundary and provenance}
The theorem applies to the supplied exact table. It does not validate floating-point
values or external coverage. Lean checked the exact statement; the dossier and proof were
AI-assisted.
\end{document}
